\section{Introduction}
\label{sec:introduction}

A camera can provide input for an interface through simple hand gestures. In
a live system, however, the hand must first be visible and detected, and the
output must remain useful as the user changes pose. A classifier can score
highly on cropped images with valid landmarks while the live system still
misses gestures because of framing, occlusion, or a short hold. We study these
stages separately in a prototype that turns recognized digits into speech or
music.

The task is to recognize a static, one-hand pose corresponding to a digit
from 0 to 9, or to identify it as belonging to a small set of non-target
poses. The digit labels reflect the project's choice of hand shapes. The
primary research question is whether explicit geometry, including finger
bending and palm orientation, provides useful information beyond normalized
landmark coordinates. To test this, we train two lightweight classifiers with
three input representations on the same split. We also audit the hand
detector and apply the browser's confidence rule to the retained still images
to show where the reported classification accuracy applies.

The implementation uses the pretrained MediaPipe hand tracker
\cite{zhang2020mediapipe,mediapipehandlandmarker}. The project implements the
hand representation, classifier training and comparison, and the browser
decision and feedback logic. Code and experiment instructions are available in
the \href{https://github.com/centaurus-xie/csit5410-hand-gesture-recognition}
{project repository}.
